\documentclass[UTF8,a4paper,11pt,openany]{ctexbook}
\usepackage{../common/statlearnbook}
\SLSetBookMetadata{第 5 册}{统计学习方法初学者讲义 v2.7.0}{采样方法、主题模型与图排序}
\begin{document}
\setcounter{page}{714}
\setcounter{chapter}{36}
\setcounter{section}{1}
\setcounter{figure}{2}
\pagestyle{headings}
\chaptermark{\PageRank{}算法}
\sectionmark{基本 \PageRank{}}
不可约但周期为$2$的链仍可有唯一平稳分布，却可能在两个状态间振荡；
可约链则可能有多个闭类和多个平稳分布。图\ref{fig:V5-C07-periodic-dangling}
在本章列向量、列随机约定下并列展示周期振荡与悬挂零列两类故障及其不同修复顺序。
% v2.7.0 figure UID: FIG-P716-01
\tikzset{slfig-FIG-P716-01/.style={font=\fontsize{9.6pt}{11.6pt}\selectfont,every path/.append style={line cap=round,line join=round}}}
\pgfkeysifdefined{/tikz/alt}{}{\tikzset{alt/.code={}}}
\begin{figure}[htbp]
\centering
\begin{tikzpicture}[slfig-FIG-P716-01,
  >=Stealth,
  every node/.style={font=\fontsize{9.6pt}{11.6pt}\selectfont,align=center},
  panel/.style={draw=SLRuleGray,line width=0.72pt,rounded corners=3pt,
    minimum width=76mm,minimum height=90mm,inner sep=0pt},
  title left/.style={font=\fontsize{10.2pt}{12.2pt}\selectfont\bfseries,
    text=SLAccentOrange,inner sep=0pt},
  title right/.style={font=\fontsize{10.2pt}{12.2pt}\selectfont\bfseries,
    text=SLDeepBlue,inner sep=0pt},
  page/.style={circle,draw=SLDeepBlue,fill=SLSoftBlue,line width=0.82pt,
    minimum size=10mm,inner sep=1pt},
  dangling/.style={rectangle,rounded corners=2pt,draw=SLAccentOrange,
    fill=white,line width=1.0pt,minimum width=12mm,minimum height=10mm,inner sep=1pt},
  link/.style={->,draw=SLMidBlue,line width=0.94pt},
  transition label/.style={font=\fontsize{9.6pt}{11.6pt}\selectfont,inner sep=0pt},
  cell/.style={draw=SLRuleGray,line width=0.58pt,minimum size=6.4mm,inner sep=0pt},
  state one/.style={draw=SLDeepBlue,solid,fill=SLSoftBlue,line width=0.76pt,
    rounded corners=2pt,minimum width=13mm,minimum height=11mm,inner sep=1pt},
  state two/.style={draw=SLAccentOrange,dashed,fill=white,line width=0.82pt,
    rectangle,minimum width=13mm,minimum height=11mm,inner sep=1pt},
  sequence/.style={->,draw=SLTextGray,densely dotted,line width=0.74pt},
  fact/.style={draw=SLDeepBlue,solid,fill=SLSoftBlue,line width=0.76pt,
    rounded corners=2pt,text width=66mm,minimum height=9mm,
    inner xsep=1.6mm,inner ysep=1.2mm},
  warning/.style={draw=SLAccentOrange,dashed,fill=white,line width=0.82pt,
    rounded corners=2pt,text width=66mm,minimum height=12mm,
    inner xsep=1.6mm,inner ysep=1.2mm},
  repair/.style={draw=SLDeepBlue,dash pattern=on 4.5pt off 2pt on 1pt off 2pt,
    fill=SLWarmGray,line width=0.78pt,rounded corners=2pt,text width=66mm,
    inner xsep=1.6mm,inner ysep=1.2mm},
  >={Stealth[length=2.15mm,width=1.55mm]},
  alt={对象—关系—结论：左面板保持A与B双向确定转移的二周期拓扑，列随机矩阵为零一交换矩阵；从第一基向量出发，时刻零、二回到第一基向量，时刻一落在第二基向量。该链存在唯一平稳分布二分之一乘全一向量，但从此初值的逐步序列不收敛，修复方向只概括为阻尼或惰性。右面板保持C指向零出度D的拓扑，矩阵第二列和为零且作用于第二基向量得到零向量，所以它不是列随机概率核并会把质量一直接丢到零。一般PageRank必须先用概率向量v回填悬挂列形成列随机S，再构造阻尼矩阵G。颜色之外，圆形与矩形、实线与虚线及点划线共同编码状态和结论。}]

\node[panel] at (-4.05,0) {};
\node[title left] at (-4.05,4.00) {故障一：周期$2$振荡};
\node[page] (a) at (-5.55,2.85) {A};
\node[page] (b) at (-2.55,2.85) {B};
\draw[link] (a) to[bend left=17] node[transition label,above=1.6mm] {$1$} (b);
\draw[link] (b) to[bend left=17] node[transition label,below=1.8mm] {$1$} (a);
\node at (-4.05,1.65)
  {$M_{\mathrm p}=\begin{bmatrix}0&1\\1&0\end{bmatrix}$，
   $\boldsymbol r^{(0)}=\boldsymbol e_1$};

\node[state one] (s0) at (-5.65,0.45) {$t=0$\\$\boldsymbol e_1$};
\node[state two] (s1) at (-4.05,0.45) {$t=1$\\$\boldsymbol e_2$};
\node[state one] (s2) at (-2.45,0.45) {$t=2$\\$\boldsymbol e_1$};
\draw[sequence] (s0) -- (s1);
\draw[sequence] (s1) -- (s2);

\node[fact] at (-4.05,-0.80)
  {\textbf{平稳分布：存在且唯一}\quad
   $\boldsymbol r_\star=(\tfrac12,\tfrac12)^{\mathsf T}$};
\node[warning] at (-4.05,-2.10)
  {\textbf{从该初值逐步收敛：无}\\
   $M_{\mathrm p}^{2q}\boldsymbol e_1=\boldsymbol e_1$，
   $M_{\mathrm p}^{2q+1}\boldsymbol e_1=\boldsymbol e_2$};
\node[repair,minimum height=10mm] at (-4.05,-3.55)
  {\textbf{修复方向（概括）：}阻尼／惰性};

\node[panel] at (4.05,0) {};
\node[title right] at (4.05,4.00) {故障二：悬挂零列};
\node[page] (c) at (2.65,2.85) {C};
\node[dangling] (d) at (5.45,2.85) {D};
\draw[link] (c) -- node[transition label,above=1.6mm] {$1$} (d);
\node[text=SLAccentOrange,inner sep=0pt] at (6.90,2.85) {零出度};
\node[inner sep=0pt] at (2.55,1.25) {$M_{\mathrm d}=$};
\matrix (md) [matrix of nodes,nodes={cell},row sep=-\pgflinewidth,column sep=-\pgflinewidth]
  at (4.05,1.25) {$0$&$0$\\$1$&$0$\\};
\node[draw=SLAccentOrange,dashed,line width=0.86pt,
  fit=(md-1-2)(md-2-2),inner sep=0pt] {};
\node[align=left,text=SLTextGray,inner sep=0pt] at (6.15,1.25)
  {第$2$列和$=0$};

\node[fact,fill=white] at (4.05,-0.10)
  {$M_{\mathrm d}\boldsymbol e_2=\boldsymbol0$，\quad
   $\boldsymbol1^{\mathsf T}M_{\mathrm d}=(1,0)\ne\boldsymbol1^{\mathsf T}$};
\node[warning] at (4.05,-1.45)
  {\textbf{不是列随机概率核，并丢失质量}\\
   $\boldsymbol r^{(0)}=\boldsymbol e_2\Rightarrow
   \boldsymbol r^{(1)}=\boldsymbol0$，总质量$1\to0$};
\node[repair,minimum height=19mm] at (4.05,-3.25)
  {\textbf{顺序不可交换：}\\
   先用$\boldsymbol v$回填悬挂列：
   $S=M_{\mathrm d}+\boldsymbol v\boldsymbol e_2^{\mathsf T}$（列随机）；\\
   再阻尼：$G=dS+(1-d)\boldsymbol v\boldsymbol1^{\mathsf T}$，$0\le d<1$。};
\end{tikzpicture}
\captionsetup{width=.94\linewidth}
\caption{基本PageRank的两类结构性故障：周期$2$链可有唯一平稳分布却从给定初值振荡；悬挂零列不是列随机核并会丢失质量，须先回填再阻尼。}
\label{fig:V5-C07-periodic-dangling}
\end{figure}

\FloatBarrier
\noindent\textbf{读图检查。}左栏先核对
$M_{\mathrm p}\boldsymbol e_1=\boldsymbol e_2$、
$M_{\mathrm p}\boldsymbol e_2=\boldsymbol e_1$，再区分“唯一平稳分布
$(1/2,1/2)^{\mathsf T}$存在”与“从$\boldsymbol e_1$逐步收敛”（后者不成立）；
右栏核对$M_{\mathrm d}$第$2$列和为$0$且
$M_{\mathrm d}\boldsymbol e_2=\boldsymbol0$，故它不是列随机概率核并丢失质量，
最后复述一般\PageRank{}必须先以$\boldsymbol v$回填悬挂列得到$S$，再由$S$构造阻尼$G$。
\end{document}
